

\chapter{Building a Character}


\section{1. Race/Species}
A character's \textit{Race} is the species that serves as the 
statistical foundation for that character.

\vspace{0.4em}

When you create your character, first choose a race.

% TODO Custom Race/Species box
\begin{Example}
  \textbf{HUMAN (Default)}\\

  % TODO Size: Normal, Move Distance: 3"

  \textbf{Base Reservoirs:}
  \vspace{-.4em}
  \begin{Table}{l l}
    Life & 20 \\
    Mana & 10 \\
  \end{Table}
  \vspace{-.5em}

  \textbf{Racial Traits:}
  \begin{itemize}
    \item +1 to any Attribute (chosen after rolling Attributes)
    \item +1 to any Attribute (chosen after rolling Attributes)
  \end{itemize}

  \textbf{Racial Modifiers:}
  \vspace{-.5em}
  \begin{Table}{X c}
    \TableHeader{Attribute} & \TableHeader{Modifier} \\
    W, I, Z, A, R, D, S & +0
  \end{Table}
  \vspace{-.5em}
\end{Example}


\subsection{Racial Traits \& Modifiers}
\textit{Racial Traits} are the perks and abilities available 
to a character upon creation.

\vspace{0.4em}

\textit{Racial Modifiers} should be used as the default attribute modifiers 
for the GM's characters of that race.

\vspace{0.4em}

When creating a player character, 
add the Racial Modifiers to the W.I.Z.A.R.D.S. attributes as they are rolled.


\subsection{Hybrid Races}
At the GM's discretion, \textit{Hybrid} characters may be created 
from a combination of two or more races.
These races become the character's \textit{Ancestral Races}.

\vspace{0.4em}

To create a character with a hybrid race:
\begin{itemize}
  \item Use the average of the Ancestral Race's Life and Mana reservoirs, 
    rounding down.
  \item Use the average of the Ancestral Races' Attribute Modifiers, 
    rounding down.
  \item Use the Creature Size that is closest to Normal.
  \item Start without access to racial traits.
\end{itemize}

Whenever a character with a hybrid race is able to level-up, 
they may instead "unlock" a racial trait from one of their Ancestral Races.



\section{2. W.I.Z.A.R.D.S. Attributes}
Each character has a set of \textit{Attributes} to represent their overall capability,
which can be remembered using the acronym W.I.Z.A.R.D.S.

\begin{Table}{l X}
  \TableHeader{Attribute} & \TableHeader{Traits Covered} \\
  Wisdom       & Intuition, Insight, Depth of Knowledge \\
  Intelligence & Complex Analysis, Abstract Reasoning \\
  Zest         & Charisma, Flair, Force of Personality \\
  Agility      & Mobility, Reactivity, Balance \\
  Resilience   & Stamina, Endurance, Natural Immunity \\
  Dexterity    & Precision, Fine Motor Control \\
  Strength     & Sturdiness, Raw Muscular Force \\
\end{Table}

W.I.Z.A.R.D.S. Attributes are used to limit the action abilities, 
passive perks, and equipment items available to the character,
representing the character's practical capacity to use the ability or item.

\vspace{0.4em}

If a character's Attribute scores are lower than the values required by
an ability, perk, or item, 
none of the effects it grants may be activated or applied.

\vspace{0.4em}

When you create your character, roll \Roll{1d10} for each Attribute, 
then add the character's Racial Modifiers to each of the results.
These become your character's starting Attributes.

\subsection{Value Range}
Attribute values can range from \textbf{0 - 10} by default. 
Modifiers to an attribute value (e.g., from perks or equipment)
cannot increase it above the maximum or decrease it below the minimum 
of this range.

\subsection{Attribute Modifiers}
Attributes also determine the default modifiers a character uses for 
Story Rolls.

%\vspace{0.4em}
%
%These may be referred to as \textit{W.I.Z.A.R.D.S. Modifiers}, 
%or individually as the character's \textit{<Attribute> Modifier}
%(e.g., ``Zest Modifier").

\begin{Table}{c c}
  \TableHeader{Value} & \TableHeader{Modifier} \\
  0  & -5 \\
  1  & -4 \\
  2  & -3 \\
  3  & -2 \\
  4  & -1 \\
  5  & +0 \\
  6  & +1 \\
  7  & +2 \\
  8  & +3 \\
  9  & +4 \\
  10 & +5 \\
\end{Table}



\section{3. Reservoirs}
Character \textit{Reservoirs} are resource pools that can be 
reduced by damage, increased by restoration, 
and expended as a cost to use abilities.


\subsection{Life}
\textit{Life} is a reservoir that represents a character's vitality 
and their ability to remain active in the game.

\vspace{0.4em}

When a character's Life is reduced to \textbf{0}, 
they perish and are removed from the game.

\vspace{0.4em}

A character's starting Life is equal to their racial Life 
plus their \textbf{Resilience} and \textbf{Strength} modifiers. 


\subsection{Mana}
\textit{Mana} is a reservoir that represents a character's magical energy
and their ability to cast spells.

\vspace{0.4em}

When a character's Mana is reduced to \textbf{0},
they may no longer cast spells, even if the Mana Cost of a spell is 0.

\vspace{0.4em}

A character's starting Mana is equal to their racial Mana 
plus their \textbf{Wisdom} and \textbf{Zest} modifiers.


\section{4. Level-Up}
A character may \textit{Level-Up} as it is created, 
and again at the GM's discretion, 
up to a maximum level of \textbf{100}.

\vspace{0.4em}

To level-up, a character may choose one:
\begin{itemize}
  \item Add a new passive perk. (See Chapter 5: Perks)
  \item Add a new action ability. (See Chapter 6: Abilities)
  \item Add a new basic spell. (See Chapter 7: Spells)
  \item Add 1 to their maximum Life and Mana.
  \item Add 1 to any Attribute.
\end{itemize}

\begin{DesignerNote}
  The pace of gaining levels will be unique to each campaign, 
  but should give players time to familiarize themselves with new abilities 
  while also providing consistent advancement.

  \vspace{0.4em}

  A good approach might be to have characters level-up during ``downtime" 
  between missions.
\end{DesignerNote}


\subsection{Starting Level}
When creating a character, determine the starting level for the campaign, 
then level-up the character that many times.

\begin{Table}{l X}
  \TableHeader{Lv.} & \TableHeader{Player} \\
    0   & New to TTRPG \\
    5   & Migrating from another TTRPG \\
    10+ & Previous experience with W.I.Z.A.R.D.S. \\
\end{Table}


%\subsection{Replacing a character}
%If a player needs to replace a character that was killed or otherwise 
%removed from the campaign,
%they may instead start at the level of their previous character.

%\begin{CruelMode}
%  Ignore this sub-rule for a more ``hardcore" experience.
%\end{CruelMode}


%\subsection{Downtime Activities}
%Some skills allow you to perform an activity if you level-up 
%in a [town](/rules/world/towns.md) that contains a certain type of Amenity.
%You may choose *one* of those activities to perform as a downtime activity.



\section{5. Armor/Apparel}
\textit{Armor} refers broadly to all apparel that can be worn by
a character to provide them with protective benefits.

\vspace{0.4em}

Each piece of Armor has an \textit{Armor Type}:
\begin{itemize}
  \item Head 
  \item Chest 
  \item Hands 
  \item Legs 
  \item Feet
\end{itemize}

Each piece of Armor has an \textit{Armor Class}:
\begin{itemize}
  \item Unarmored 
  \item Light 
  \item Heavy
\end{itemize}

\vspace{0.4em}

A character can have only one of each Armor Type equipped at a time.

\vspace{0.4em}

When you create your character, choose an Armor Class for each Armor Type.


\subsection{Overall Armor Class}
Your character's \textit{Overall Armor Class} is determined by the total set 
of all armor you are wearing.

\vspace{0.4em}

If 3+ slots are equipped with the same Armor Class, 
your Overall Armor Class becomes that Armor Class.
Otherwise, your Overall Armor Class is Light Armor.

\vspace{0.4em}

A character's Overall Armor Class has consequences for their Clash Save
and their Movement Distance:

\begin{Table}{l c c}
  \TableHeader{Class} & \TableHeader{Clash Save} & \TableHeader{Move Distance}\\
  Unarmored & -1 & +1" \\
  Light     & +0 & +0" \\
  Heavy     & +1 & -1" \\
\end{Table}


\subsection{Attribute Penalty}
Each piece of armor boosts or penalizes a character's Attributes 
according to the Armor Type and Class:

\begin{Table}{l c}
  \TableHeader{Type} & \TableHeader{Attribute} \\
  Head  & Wisdom, Intelligence, Zest \\
  Chest & Strength \\
  Hands & Dexterity \\
  Legs  & Resilience \\
  Feet  & Agility \\
\end{Table}

The boost/penalty is applied for each slot based on the Class of that armor:

\begin{Table}{l c}
  \TableHeader{Class} & \TableHeader{Boost/Penalty} \\
  Unarmored & +1 \\
  Light     & +0 \\
  Heavy     & -1 \\
\end{Table}

Ensure the Armor that your character equips does not reduce their Attributes 
below the requirements for the perks, abilities, and equipment
you want them to use.



\section{6. Inventory Storage}
A character may start equipped with one type of storage.

\vspace{0.4em}

\textit{Storage} is any equipment that adds inventory space to a character.
Like all wearable equipment, 
storage can be equipped to or unequipped from the body with an Action, 
provided the character has a free hand to which they then equip the storage.

\vspace{0.4em}

Each inventory item has a \textit{Storage Class}, 
and each item of storage has a set of \textit{Storage Slots}.
Items may be stored in an open Storage Slot of equal or lesser size 
to their Storage Class.

\vspace{0.4em}

Storage may also have \textit{Attach Slots}.

Attach Slots are used to add modular storage, 
or to store items that are too bulky to qualify as `Large`,
provided the item in question has the ``Attachable" property.

\vspace{0.4em}

All weapons are by default attachable, 
to allow them to be sheathed when outside of combat.


\begin{DesignerNote}
  This system is intended to promote realism, prevent inventory bloat, 
  and to encourage players to actually use all of their items 
  instead of hoarding everything the come across.

  \vspace{0.4em}

  Forcing players to be thoughtful about their inventory helps them 
  avoid the long lists of forgettable, 
  unusable items that come from less-organized or 
  less-strict inventory systems.
\end{DesignerNote}


\begin{Example}
  \textbf{Backpack:}\\
  -1 to \textbf{Agility} and \textbf{Resilience}
  \vspace{-.5em}
  \begin{Table}{l l}
    \TableHeader{Slots} & \TableHeader{\#} \\
    Large & 1 \\
    Medium & 2 \\
    Small & 5 \\
    Attach & 4 \\
  \end{Table}
  \vspace{-.5em}
\end{Example}

\begin{Example}
  \textbf{Belt:}
  \vspace{-.5em}
  \begin{Table}{l l}
    \TableHeader{Slots} & \TableHeader{\#} \\
    Small & 2 \\
    Attach & 2 \\
  \end{Table}
  \vspace{-.5em}
\end{Example}



\section{7. Equipment}
% TODO Elaborate/introduction

\subsection{Weapons}
Each character is created with both hands equipped.

\vspace{0.4em}

Choose a weapon:
\begin{itemize}
  \item Up to 2 One-Handed weapons
  \item 1 Two-Handed weapon
  \item 1 Mana Focus
\end{itemize}


\subsection{Items}
If a character is equipped with storage, 
they may fill all storage slots with standard items
that are representative of their background.


\subsection{Currency}
At the GM's discretion, a character may start with currency 
appropriate to the campaign economy and the character's background. 

\vspace{0.4em}

The character must have storage capacity for any currency they receive.


